% vim: set spell spelllang=en_us :


\green{We first describe the structures generated by the two atom-substitution procedures 
and the reported global-minimum (GM) structures used for comparison. Their
geometries and spin states define the models used in the bonding and \ce{CO}
adsorption analysis below. Figure~\ref{fig:crossed_models} summarizes the two
routes. Starting from reported \ce{Pt10} GM\cite{ugartemendia:2022},
selected Pt atoms are replaced by Ge to obtain the 
\ce{Pt6Ge4^{der}}, \ce{Pt5Ge5^{der}}, and \ce{Pt4Ge6^{der}} derived clusters. In the reverse
\ce{Pt5Ge5} to \ce{Pt10} derivation, all Ge atoms in the reported GM
\ce{Pt5Ge5}\cite{ugartemendia:2022} structure are replaced by Pt, leading to a new \ce{Pt10^{der}}
isomer. The newly derived isomers are optimized and the changes in the geometry
and the bonding patterns are examined.

Here, \ce{Pt10^{GM}} and \ce{Pt5Ge5^{GM}} denote the reoptimized reported
global minima\cite{ugartemendia:2022}, and \ce{der} denotes a derivative structure.}


\begin{figure}[H]
\centering
\begin{tikzpicture}[
  x=1cm,
  y=1cm,
  route/.style={font=\sffamily\bfseries\large, align=center},
  note/.style={font=\sffamily\small, align=center, text=black!70},
  molname/.style={font=\sffamily\bfseries\small, align=center},
  molnote/.style={font=\sffamily\scriptsize, align=center, text=black!70},
  arrow/.style={->, >=stealth, line width=0.7pt, black!85}
]

% Image sizes.  These are intentionally defined here so the layout can be
% adjusted directly from LaTeX.
\def\parentw{0.19\linewidth}
\def\parentwide{0.17\linewidth}
\def\molw{0.15\linewidth}
\def\molwb{0.15\linewidth}

% Molecule images.  Draw these first because the PNGs have white backgrounds.
\node[inner sep=0pt] (pt10parent) at (-3.30,2.75)
  {\includegraphics[width=\parentw]{pt10.png}};
\node[inner sep=0pt] (pt5ge5parent) at (3.45,2.75)
  {\includegraphics[width=\parentwide]{pt5ge5_orig.png}};

\node[inner sep=0pt] (pt6ge4) at (-5.15,-0.10)
  {\includegraphics[width=\molw]{pt6ge4.png}};
\node[inner sep=0pt] (pt5ge5) at (-3.25,-0.10)
  {\includegraphics[width=\molw]{Pt5Ge5_fromPt10Parent.png}};
\node[inner sep=0pt] (pt4ge6) at (-1.15,-0.10)
  {\includegraphics[width=\molw]{pt4ge6.png}};
\node[inner sep=0pt] (pt10frompt5ge5) at (3.45,-0.10)
  {\includegraphics[width=\molwb]{pt10_from_pt5ge5.png}};

% Route arrows.
\draw[arrow] (-3.82,1.18) -- (-4.90,0.80);
\draw[arrow] (-3.30,1.18) -- (-3.25,0.82);
\draw[arrow] (-2.78,1.18) -- (-1.55,0.80);
\draw[arrow] (3.45,1.18) -- (3.45,0.85);

% Headers and labels.  Draw these last so they remain visible over the PNG
% image boxes.
\node[route] at (-3.30,4.58) {\ce{Pt10^{GM}} starting structure};
\node[note] at (-3.30,4.16) {Ge substitution followed by optimization};
\node[route] at (3.45,4.58) {\ce{Pt5Ge5^{GM}} starting structure};
\node[note] at (3.45,4.16) {Ge-to-Pt substitution followed by optimization};

\node[molname] at (-5.15,-1.18) {\ce{Pt6Ge4^{der}}};
\node[molname] at (-3.25,-1.18) {\ce{Pt5Ge5^{der}}};
\node[molname] at (-1.35,-1.18) {\ce{Pt4Ge6^{der}}};
\node[molname] at (3.45,-1.18) {\ce{Pt10^{der}}};

\node[note] at (-0.70,-2.22)
  {Upper row: reported GM structures. Lower row: optimized structures after atom substitution.};

\end{tikzpicture}

\caption{\green{Derivative structures used in the
\ce{CO}-adsorption study. The upper row shows the reoptimized reported GM
structures, and the lower row shows the optimized derivatives obtained after
Pt/Ge substitution. Labels identify the compositions and derivation direction.}}
\label{fig:crossed_models}
\end{figure}

\subsection{Starting-structure effects in \ce{Pt10} and Pt--Ge clusters}

Figure~\ref{fig:crossed_models} compares the previously reported Pt10 and Pt5Ge5 global-minimum (GM) structures, from-now on \ce{Pt_10^{GM}} and \ce{Pt5Ge5^{GM}}, which will serve as our reference. On the one hand, the monometallic \ce{Pt10^{GM}} structure is a compact, highly symmetric pyramidal
cluster with a nonet ground state. \blue{The equimolar
\ce{Pt5Ge5^{GM}} structure is also compact, but it has a near-fourfold
cage geometry. Four Pt atoms form a square-like group capped by Ge,
whereas four Ge atoms form the opposite group capped by Pt. It has a
singlet ground state and chemically distinct Pt and Ge sites.} The two starting
structures therefore differ in both geometry, spin-state, and chemical bonding before any atom substitution is made.

In order to understand what is the interplay between the chemical bond (metallic/delocalized vs covalent/localized) and the geometry, we consider these \ce{Pt10^{GM}} and \ce{Pt5Ge5^{GM}} models and play with the composition to evaluate the dependence bwteen the geometry, chemical bond and composition and the impact of these factors in the CO-Pt interaction strength. Replacing selected Pt atoms in \ce{Pt10^{GM}} with Ge gives the
\ce{Pt6Ge4^{der}}, \ce{Pt5Ge5^{der}}, and \ce{Pt4Ge6^{der}} models in the lower left part of
Figure~\ref{fig:crossed_models}. The replacement is done with the aim of obtaining alternating Pt and Ge atoms, trying to mimick the  \ce{Pt5Ge5^{GM}} structure.  The Ge
substitution and subsequent relaxation lowers the symmetry and creates Pt sites
with different Pt/Ge environments, unlike in \ce{Pt5Ge5^{GM}} . The ground spin-state is singlet for all the newly obtained PtGe clusters. Therefore, the preference for closed-shell electronic structures in PtGe alloys retains even with symmetries that are not preferred by such compositions. In the reverse derivation, replacing all Ge
atoms in \ce{Pt5Ge5^{GM}} with Pt gives \ce{Pt10^{der}}, a lower-symmetry triplet structure. 

For the derivative adsorption series, this triplet
\ce{Pt10^{der}} structure is used as the pure-Pt model; the
quintet, only 0.04 eV higher in energy, is included to test the effect of spin state. The
\ce{Pt6Ge4^{der}}, \ce{Pt5Ge5^{der}}, and \ce{Pt4Ge6^{der}} structures derived from \ce{Pt10^{GM}}
are singlets. GM structures provide a second
set of adsorption calculations: \ce{Pt10^{GM}} is a nonet, \ce{Pt6Ge4^{GM}}
is a triplet, and \ce{Pt5Ge5^{GM}} and \ce{Pt4Ge6^{GM}} are singlets.
\textbf{\green{estaia bien saber  como de lejos en E estan estas nuevas estructuras respecto a los GMs.}}

\textbf{Revisar: The selected triplet wavefunction and the seven distinct triplet
\ce{Pt10CO} minima were confirmed to be orbital-stable.
The \(\ce{Pt_xGe_{10-x}^{der}}\) series  favor
the closed-shell singlet over the triplet by 0.18--0.26 eV. Thus, within this
set of models, the Ge-containing clusters have a clearer preference for
the low-spin state than \ce{Pt10^{der}}.}

The \ce{Pt10^{GM}} and \ce{Pt10^{der}} structures have the same composition
but \ce{Pt10^{GM}} is a nonet, whereas \ce{Pt10^{der}} is a triplet. At the
present level, the latter lies 1.29 eV above the reported GM in $E_0$. This
difference combines the effects of their distinct structures and multiplicities.

\textbf{Revisar. Quiza introducir un poco que es ECoN}
Their local Pt environments also differ. The ECoN values for \ce{Pt10^{GM}}
fall into two groups: six Pt atoms have values close to 5.68, while four have
an ECoN of 3.00. In the \ce{Pt10^{der}} triplet, the ECoN values span
3.03--6.01 and define five pairs of similar Pt environments. All ECoN values
for the individual Pt atoms are given in Supporting
Table~\ref{tab:tfg_pt10_bare_econ}. These values show that the is a homogeneous, compact, and symmetric structure compared to the derivative.

We made the same comparison for the Pt--Ge compositions. The total ECoN
values are different for the GM and derivative structures. For \ce{Pt6Ge4},
the GM structure has a narrower range (3.84--4.87) than the derivative
structure (2.77--4.60). The difference is also clear for \ce{Pt5Ge5}, where
the derivative structure has a wider ECoN range than the GM structure. The
two \ce{Pt4Ge6} structures also show different distributions of total ECoN.

The initial Ge contribution, $x_{\rm Ge}^{0}$, gives a second comparison.
This quantity is the fraction of the effective coordination of the Pt site
provided by Ge neighbors before \ce{CO} adsorption.
Further details of the definition and the partitioning of the ECoN are given
in the Supporting Information (Section~\ref{sec:supporting_econ_definition}).
The Pt sites in \ce{Pt5Ge5^{GM}} are all Ge-rich
($x_{\rm Ge}^{0}=0.930$--0.998), whereas the derivative structure includes a
site with a lower Ge contribution ($x_{\rm Ge}^{0}=0.745$). The two
\ce{Pt4Ge6} structures are Ge-rich in both cases, while the \ce{Pt6Ge4}
structures contain a broader range of local Ge environments. The descriptors
are given in Table~\ref{tab:tfg_crossed_ptge_econ_descriptors} and Supporting
Table~\ref{tab:tfg_gm_ptge_econ_pre_descriptors}. Thus, changing the starting
structure changes both the coordination and the local Pt/Ge environment
available for \ce{CO} adsorption, even at the same composition. The next
section examines the corresponding bonding patterns.
